% Prose rendering of PrimeTensor/Completion.lean.
% Fragment: no \documentclass, no preamble, no document environment.

\section{Intrinsic multiplicative completion}
\label{Completion:sec}

\leanfile{PrimeTensor/Completion.lean}

\subsection*{What this file establishes}

Everything so far has been about a single pair of magnitudes: how close two
elements of $\name{MulRat}$ are, measured on the ladder of
\texttt{PrimeTensor/Scale.lean} and composed by the law of
\texttt{PrimeTensor/Scale/Laws.lean}.  This file passes to sequences.  It
defines what it means for a stream of magnitudes to be Cauchy for that ladder,
proves that eventual agreement of two such streams is an equivalence relation,
and forms the quotient
\[
  \name{MulReal} \defeq
    \{\text{Cauchy streams}\} \,/\, \{\text{asymptotic equivalence}\},
\]
the carrier the rest of the development calls the multiplicative reals.

Two things have to be built before the analysis can start, and they occupy the
first half of the file.  A Cauchy condition speaks of a \emph{tail} of the
index set, so the depths need an order; $\name{Depth}$ of
\texttt{PrimeTensor/Depth.lean} carries none, and one is defined here as an
inductive relation.  And combining two Cauchy conditions needs a common tail,
so the depths need a binary upper bound; that is $\name{join}$, the analogue of
$\max(N_1, N_2)$.

The organising decision inherited from the whole development is that the
\emph{tolerance index and the stage index are the same type}.  A classical
Cauchy condition reads ``for every $\varepsilon > 0$ there is an $N$'': two
different kinds of object, one a magnitude and one a stage.  Here it reads
``for every level $\ell$ there is an anchor $a$'', and both are depths.  Nothing
is quantified over that is not already in the object language, and in
particular no $\varepsilon$ has to be halved --- the halving of the classical
proof becomes $\ctor{succ}$, which
\texttt{PrimeTensor/Scale/Laws.lean} arranged to be exactly right.

\subsection{An order on the depths}

\begin{definition}[\lean{Depth.AtOrAfter}]\label{Completion:def:atorafter}
$\name{AtOrAfter}$ is the inductive relation on $\name{Depth}$ generated by
\[
  \ctor{here} : \name{AtOrAfter}(n, n),
  \qquad
  \ctor{later} : \name{AtOrAfter}(a, n) \to
    \name{AtOrAfter}(a, \ctor{succ}\,n).
\]
Read $\name{AtOrAfter}(a, n)$ as ``$n$ lies in the tail beginning at $a$'',
that is, $\size{a} \leq \size{n}$.
\end{definition}

\begin{leancode}
inductive AtOrAfter : Depth → Depth → Prop where
  | here (n : Depth) : AtOrAfter n n
  | later {anchor n : Depth} :
      AtOrAfter anchor n → AtOrAfter anchor (.succ n)
\end{leancode}

The recursion is on the \emph{right} endpoint: a proof of
$\name{AtOrAfter}(a,n)$ is literally the number of successor steps from $a$ up
to $n$, packaged as $\size{n} - \size{a}$ applications of $\ctor{later}$ to a
single $\ctor{here}$.  Induction on such a proof is therefore induction on
that difference, which is what the three lemmas below use.

\begin{lemma}[\lean{Depth.one\_atOrAfter}]\label{Completion:lem:one-atorafter}
$\name{AtOrAfter}(\ctor{one}, n)$ for every $n$.
\end{lemma}

\begin{leancode}
theorem one_atOrAfter : ∀ n : Depth, AtOrAfter .one n
  | .one => .here .one
  | .succ n => .later (one_atOrAfter n)
\end{leancode}

\begin{proof}
By structural recursion on $n$.  At $\ctor{one}$ the two endpoints coincide, so
$\ctor{here}$ applies; at $\ctor{succ}\,n$ the recursive call supplies
$\name{AtOrAfter}(\ctor{one}, n)$ and $\ctor{later}$ raises it.
\end{proof}

So $\ctor{one}$ is a least element, and the whole index set is one tail --- the
tail anchored at $\ctor{one}$.  This is what lets the constant stream below
choose $\ctor{one}$ as its anchor and mean ``everywhere''.

\begin{lemma}[\lean{Depth.succ\_atOrAfter}]\label{Completion:lem:succ-atorafter}
$\name{AtOrAfter}(a,b)$ implies
$\name{AtOrAfter}(\ctor{succ}\,a, \ctor{succ}\,b)$.
\end{lemma}

\begin{proof}
By induction on the given proof.  In the $\ctor{here}$ case both endpoints are
raised to the same successor, so $\ctor{here}$ applies again.  In the
$\ctor{later}$ case the induction hypothesis gives
$\name{AtOrAfter}(\ctor{succ}\,a, \ctor{succ}\,n)$ and one more $\ctor{later}$
reaches $\ctor{succ}(\ctor{succ}\,n)$.
\end{proof}

\begin{lemma}[\lean{Depth.atOrAfter\_trans}]\label{Completion:lem:trans}
$\name{AtOrAfter}(a,b)$ and $\name{AtOrAfter}(b,c)$ imply
$\name{AtOrAfter}(a,c)$.
\end{lemma}

\begin{proof}
By induction on the second proof, the first being carried along unchanged.  The
$\ctor{here}$ case is $b = c$, where the first proof \emph{is} the conclusion;
the $\ctor{later}$ case appends one more step to the induction hypothesis.  In
effect the two step-counts are added.
\end{proof}

\begin{remark}[Reflexive and transitive, and nothing more is proved]
\label{Completion:rem:preorder}
$\ctor{here}$ gives reflexivity and Lemma~\ref{Completion:lem:trans} gives
transitivity, so $\name{AtOrAfter}$ is a preorder; it is in fact the usual
$\leq$ transported along $d \mapsto \size{d}$, hence also antisymmetric and
total.  Neither of those is proved, and neither is needed: a Cauchy argument
only ever needs to know that a tail is nonempty, that tails shrink, and that
two tails meet.  Nor is $\name{AtOrAfter}$ registered as a $\ctor{Preorder}$
instance, so no order notation and no library lemma about $\leq$ is in scope
--- the three lemmas above are the entire theory of it available downstream.
\end{remark}

\subsection{Common tails}

\begin{definition}[\lean{Depth.join}]\label{Completion:def:join}
$\name{join} : \name{Depth} \to \name{Depth} \to \name{Depth}$ is defined by
recursion on both arguments:
\[
  \name{join}(\ctor{one}, b) = b,
  \quad
  \name{join}(\ctor{succ}\,a, \ctor{one}) = \ctor{succ}\,a,
  \quad
  \name{join}(\ctor{succ}\,a, \ctor{succ}\,b) =
    \ctor{succ}\,\name{join}(a,b).
\]
\end{definition}

\begin{leancode}
def join : Depth → Depth → Depth
  | .one, b => b
  | .succ a, .one => .succ a
  | .succ a, .succ b => .succ (join a b)
\end{leancode}

Under $d \mapsto \size{d}$ this is $\max$, computed by peeling one successor
from each argument until one of them bottoms out.

\begin{lemma}[\lean{Depth.left\_atOrAfter} and \lean{right\_atOrAfter}]
\label{Completion:lem:join-bounds}
$\name{AtOrAfter}(a, \name{join}(a,b))$ and
$\name{AtOrAfter}(b, \name{join}(a,b))$ for all $a, b$.
\end{lemma}

\begin{leancode}
theorem left_atOrAfter : ∀ a b : Depth, AtOrAfter a (join a b)
  | .one, b => one_atOrAfter b
  | .succ a, .one => .here (.succ a)
  | .succ a, .succ b => succ_atOrAfter (left_atOrAfter a b)
\end{leancode}

\begin{proof}
Both are structural recursions following the three cases of
Definition~\ref{Completion:def:join}.  For the left bound: when $a$ is
$\ctor{one}$ the join is $b$ and Lemma~\ref{Completion:lem:one-atorafter}
applies; when $b$ is $\ctor{one}$ the join is $a$ itself and $\ctor{here}$
applies; in the successor case the recursive call is raised by
Lemma~\ref{Completion:lem:succ-atorafter}.  The right bound is the same
recursion with the two degenerate cases exchanged --- the join of $\ctor{one}$
with $b$ is $b$, so $\ctor{here}$ serves, and the join of $\ctor{succ}\,a$ with
$\ctor{one}$ is $\ctor{succ}\,a$, reached by
Lemma~\ref{Completion:lem:one-atorafter}.
\end{proof}

\begin{remark}[An upper bound, not the least one]\label{Completion:rem:join}
Only the two bounds are proved.  $\name{join}$ is not shown commutative,
associative or idempotent, and it is not shown to be a least upper bound; it
is used exclusively through Lemma~\ref{Completion:lem:join-bounds}, where all
that matters is that some single anchor lies in both tails.  Anything stronger
would be true and unused.
\end{remark}

\subsection{Cauchy streams}

\begin{definition}[\lean{MulStream} and \lean{IsMulCauchy}]
\label{Completion:def:cauchy}
A \emph{stream} is a function $s : \name{Depth} \to \name{MulRat}$, and $s$ is
\emph{multiplicatively Cauchy} when
\[
  \forall\, \ell\; \exists\, a\; \forall\, m,n :
    \name{AtOrAfter}(a,m) \wedge \name{AtOrAfter}(a,n)
    \;\longrightarrow\;
    \name{ScaleWithin}(\ell, s_m, s_n).
\]
\end{definition}

\begin{leancode}
def IsMulCauchy (s : MulStream) : Prop :=
  ∀ level : Depth,
    ∃ anchor : Depth,
      ∀ m n : Depth,
        Depth.AtOrAfter anchor m →
        Depth.AtOrAfter anchor n →
        MulRat.ScaleWithin level (s m) (s n)
\end{leancode}

\begin{designnote}[One type for both indices]\label{Completion:note:one-type}
The classical condition quantifies over a tolerance $\varepsilon > 0$ drawn
from the codomain and produces a stage $N$ drawn from the index set.  Here both
the level $\ell$ and the anchor $a$ are depths, and the codomain is not
quantified over at all.  Three consequences follow, and all three are used
below.

The stream is indexed from $\ctor{one}$, so there is no stage zero and no
convention about whether the first term counts; a tail is never empty, since
Lemma~\ref{Completion:lem:one-atorafter} makes the whole index set the tail at
$\ctor{one}$.

The levels are not quantified over an ordered field, so the classical move of
replacing $\varepsilon$ by $\varepsilon/2$ is unavailable --- and unnecessary.
Its replacement is $\ell \mapsto \ctor{succ}\,\ell$, which exists because
$\name{Depth}$ is closed under successors, and which does the same work because
\texttt{PrimeTensor/Scale/Laws.lean} proved that two steps at
$\ctor{succ}\,\ell$ compose into one at $\ell$.  Halving is thereby replaced by
a constructor.

Finally, closeness is $\name{ScaleWithin}$, which is symmetric but not
transitive at a fixed level, so the definition cannot be simplified to
consecutive terms.  The condition genuinely quantifies over both $m$ and $n$ in
the tail.
\end{designnote}

\begin{definition}[\lean{MulCauchyStream}]\label{Completion:def:stream}
$\name{MulCauchyStream}$ is the structure bundling a stream $\name{term}$ with
a proof $\name{cauchy}$ that it satisfies
Definition~\ref{Completion:def:cauchy}.
\end{definition}

\begin{lemma}[\lean{MulCauchyStream.constant} and \lean{constant\_term}]
\label{Completion:lem:constant}
For $q : \name{MulRat}$ the stream constantly equal to $q$ is Cauchy, and its
$n$-th term is $q$ definitionally.
\end{lemma}

\begin{leancode}
def constant (q : MulRat) : MulCauchyStream where
  term := fun _ => q
  cauchy := by
    intro level
    refine ⟨.one, ?_⟩
    intro m n hm hn
    exact MulRat.scaleWithin_refl level q
\end{leancode}

\begin{proof}
Given a level, take the anchor to be $\ctor{one}$, so that the tail is
everything.  Both terms are $q$, and
$\name{scaleWithin\_refl}$ of \texttt{PrimeTensor/Scale.lean} applies at every
level.  The hypotheses $\ctor{hm}$ and $\ctor{hn}$ are discarded unused.
\end{proof}

That reflexivity is unconditional matters here: had the ladder kept the
tolerance parameter of \texttt{PrimeTensor/Order.lean}, this construction would
have had to carry an admissibility hypothesis $1 < \delta$ into every constant
stream, and from there into $\name{ofRat}$ and into every rational literal of
the completion.

\subsection{Asymptotic equivalence}

\begin{definition}[\lean{MulAsymptotic}]\label{Completion:def:asymptotic}
Cauchy streams $a$ and $b$ are \emph{asymptotic} when
\[
  \forall\, \ell\; \exists\, c\; \forall\, n :
    \name{AtOrAfter}(c,n) \longrightarrow
    \name{ScaleWithin}(\ell, a_n, b_n).
\]
\end{definition}

Note that the definition mentions only the $\name{term}$ field: the Cauchy
proofs carried by $a$ and $b$ are never used, here or in any of the three
lemmas below.  The relation would be well formed on raw streams; it is stated
on the bundled structure because it is the bundled structure that gets
quotiented.

\begin{lemma}[\lean{MulAsymptotic.refl} and \lean{MulAsymptotic.symm}]
\label{Completion:lem:refl-symm}
$\name{MulAsymptotic}$ is reflexive and symmetric.
\end{lemma}

\begin{proof}
For reflexivity take the anchor $\ctor{one}$ and apply
$\name{scaleWithin\_refl}$ termwise.  For symmetry keep the anchor supplied by
the hypothesis and apply $\name{scaleWithin\_symm}$ termwise; the level is not
changed in either case.
\end{proof}

\begin{theorem}[\lean{MulAsymptotic.trans}]\label{Completion:thm:trans}
$\name{MulAsymptotic}$ is transitive.
\end{theorem}

\begin{leancode}
@[trans] theorem trans {a b c : MulCauchyStream}
    (hab : MulAsymptotic a b) (hbc : MulAsymptotic b c) :
    MulAsymptotic a c := by
  intro level
  obtain ⟨abAnchor, habTail⟩ := hab (.succ level)
  obtain ⟨bcAnchor, hbcTail⟩ := hbc (.succ level)
  let anchor := Depth.join abAnchor bcAnchor
  refine ⟨anchor, ?_⟩
  intro n hn
  have habN : Depth.AtOrAfter abAnchor n :=
    Depth.atOrAfter_trans (Depth.left_atOrAfter abAnchor bcAnchor) hn
  have hbcN : Depth.AtOrAfter bcAnchor n :=
    Depth.atOrAfter_trans (Depth.right_atOrAfter abAnchor bcAnchor) hn
  exact MulRat.scaleWithin_comp (habTail n habN) (hbcTail n hbcN)
\end{leancode}

\begin{proof}
Fix a level $\ell$.  Instantiate \emph{both} hypotheses not at $\ell$ but at
$\ctor{succ}\,\ell$, obtaining anchors $a_{ab}$ and $a_{bc}$ beyond which each
leg is close at that finer level.  Let $c^{*} = \name{join}(a_{ab}, a_{bc})$ and
take it as the anchor for $\ell$.

For $n$ at or after $c^{*}$: by Lemma~\ref{Completion:lem:join-bounds}, $c^{*}$
is at or after each of $a_{ab}$ and $a_{bc}$, so
Lemma~\ref{Completion:lem:trans} places $n$ in both tails.  Both hypotheses
therefore apply at $n$, giving
$\name{ScaleWithin}(\ctor{succ}\,\ell, a_n, b_n)$ and
$\name{ScaleWithin}(\ctor{succ}\,\ell, b_n, c_n)$, and
$\name{scaleWithin\_comp}$ of \texttt{PrimeTensor/Scale/Laws.lean} composes
them into $\name{ScaleWithin}(\ell, a_n, c_n)$.
\end{proof}

\begin{designnote}[The $\varepsilon/2$ argument, without $\varepsilon/2$]
\label{Completion:note:halving}
The classical proof of this statement asks each hypothesis for
$\varepsilon/2$, takes $N = \max(N_1, N_2)$, and adds the two halves with the
triangle inequality.  The proof above has the same three moves and none of the
arithmetic: it asks each hypothesis for $\ctor{succ}\,\ell$, takes
$\name{join}$ of the two anchors, and combines with the composition law.

Each replacement is forced by what the object language contains.  Halving a
tolerance would need division in the tolerance space; the tolerances here are
$2^{1/2^{k}}$ and are not elements of $\name{MulRat}$ at all, as
\texttt{PrimeTensor/Scale.lean} recorded, so they cannot be divided --- but the
\emph{index} of the tolerance can be raised by one, and by
\texttt{PrimeTensor/Scale/Laws.lean} raising it by one is precisely halving.
Taking a maximum would need an order on the index set with binary suprema; the
index set is $\name{Depth}$, which carries neither until
Definition~\ref{Completion:def:atorafter} and
Definition~\ref{Completion:def:join} supply exactly the two facts used.  This
is why the first half of this file exists.

The cost is that transitivity is proved by strengthening the request, so the
hypotheses must be universally quantified over levels --- which they are, and
which is why $\name{MulAsymptotic}$ could not have been stated at one fixed
level.
\end{designnote}

\begin{definition}[\lean{MulAsymptotic.setoid}]\label{Completion:def:setoid}
The setoid on $\name{MulCauchyStream}$ whose relation is
$\name{MulAsymptotic}$ and whose proof of equivalence is the triple assembled
from Lemma~\ref{Completion:lem:refl-symm} and
Theorem~\ref{Completion:thm:trans}.
\end{definition}

\subsection{The completion}

\begin{definition}[\lean{MulReal} and its constructors]
\label{Completion:def:mulreal}
\begin{align*}
  \name{MulReal} &\defeq
    \name{Quotient}(\name{MulAsymptotic.setoid}), \\
  \name{ofStream}(s) &\defeq [s], \\
  \name{ofRat}(q) &\defeq \name{ofStream}(\name{constant}(q)), \\
  \name{one} &\defeq \name{ofRat}(1),
\end{align*}
with $\name{one}$ registered as the $\ctor{One}$ instance of
$\name{MulReal}$.
\end{definition}

\begin{leancode}
abbrev MulReal := Quotient MulAsymptotic.setoid
\end{leancode}

\begin{remark}[What the carrier is, and what it is not yet]
\label{Completion:rem:what-remains}
$\name{MulReal}$ is a type with a distinguished point and a map in from
$\name{MulRat}$, and that is the whole of its structure at the end of this
file.  There is no multiplication, no inverse, no order, and no lemma
relating $\name{ofRat}$ to the multiplication of $\name{MulRat}$; nor is
$\name{ofRat}$ shown injective, which is the statement that distinct rationals
stay distinct in the completion and needs an argument that two constant streams
asymptotic at every level are equal.  Even $\name{one}$ is so far only a point:
nothing says it is a unit, there being nothing for it to be a unit for.

The reason the file stops here is that each of those requires the same
ingredient --- a proof that the operation in question respects
$\name{MulAsymptotic}$, which is a fresh $\varepsilon/2$-style argument per
operation --- and the machinery for those arguments is exactly what has just
been assembled.
\end{remark}

\begin{remark}[The Cauchy field is inert so far]\label{Completion:rem:inert}
No proof in this file reads the $\name{cauchy}$ field of a stream: the three
equivalence lemmas use only $\name{term}$, and the constant stream supplies its
proof without anyone consuming one.  The field is nonetheless load-bearing for
what follows, since a quotient of arbitrary streams by eventual agreement is
not a completion of anything.  It first does work when an operation has to be
shown to send Cauchy streams to Cauchy streams.
\end{remark}

\subsection*{Summary of the correspondence}

\begin{center}
\small
\begin{tabular}{@{}lll@{}}
\hline
Lean declaration & Kind & Here \\
\hline
\lean{Depth.AtOrAfter}              & inductive  & Def.~\ref{Completion:def:atorafter} \\
\lean{Depth.one\_atOrAfter}         & theorem    & Lem.~\ref{Completion:lem:one-atorafter} \\
\lean{Depth.succ\_atOrAfter}        & theorem    & Lem.~\ref{Completion:lem:succ-atorafter} \\
\lean{Depth.atOrAfter\_trans}       & theorem    & Lem.~\ref{Completion:lem:trans} \\
\lean{Depth.join}                   & definition & Def.~\ref{Completion:def:join} \\
\lean{Depth.left\_atOrAfter}        & theorem    & Lem.~\ref{Completion:lem:join-bounds} \\
\lean{Depth.right\_atOrAfter}       & theorem    & Lem.~\ref{Completion:lem:join-bounds} \\
\lean{MulStream}                    & abbrev     & Def.~\ref{Completion:def:cauchy} \\
\lean{IsMulCauchy}                  & definition & Def.~\ref{Completion:def:cauchy} \\
\lean{MulCauchyStream}              & structure  & Def.~\ref{Completion:def:stream} \\
\lean{MulCauchyStream.constant}     & definition & Lem.~\ref{Completion:lem:constant} \\
\lean{MulCauchyStream.constant\_term} & simp thm & Lem.~\ref{Completion:lem:constant} \\
\lean{MulAsymptotic}                & definition & Def.~\ref{Completion:def:asymptotic} \\
\lean{MulAsymptotic.refl}           & theorem    & Lem.~\ref{Completion:lem:refl-symm} \\
\lean{MulAsymptotic.symm}           & theorem    & Lem.~\ref{Completion:lem:refl-symm} \\
\lean{MulAsymptotic.trans}          & theorem    & Thm.~\ref{Completion:thm:trans} \\
\lean{MulAsymptotic.setoid}         & definition & Def.~\ref{Completion:def:setoid} \\
\lean{MulReal}                      & abbrev     & Def.~\ref{Completion:def:mulreal} \\
\lean{MulReal.ofStream}             & definition & Def.~\ref{Completion:def:mulreal} \\
\lean{MulReal.ofRat}                & definition & Def.~\ref{Completion:def:mulreal} \\
\lean{MulReal.one}                  & definition & Def.~\ref{Completion:def:mulreal} \\
\lean{One MulReal}                  & instance   & Def.~\ref{Completion:def:mulreal} \\
\hline
\end{tabular}
\end{center}

\noindent
Every declaration of \texttt{PrimeTensor/Completion.lean} appears above,
together with the two constructors $\ctor{here}$ and $\ctor{later}$ of
$\name{AtOrAfter}$ and the two fields $\name{term}$ and $\name{cauchy}$ of
$\name{MulCauchyStream}$.  The file contains no \lean{sorry}, no axiom, and no
unproved named proposition.  Remarks~\ref{Completion:rem:preorder},
\ref{Completion:rem:join}, \ref{Completion:rem:what-remains}
and~\ref{Completion:rem:inert} are stated for the reader and are not
declarations of the file.
